\documentclass[12pt]{article}

%paquete para poder cambiar interlineado
\usepackage{setspace}
\renewcommand{\baselinestretch}{1.1} %cambio interlineado
\setlength{\parskip}{0.8em} % para q haya espacio entre parrafos

\usepackage{graphicx}                  % Para tener 'includegraphics'
%\usepackage[bidi=default]{babel}      % Porque usamos español
\usepackage{float}                     % Para que las figuras aparezcan en el lugar donde las ponemos 
% \usepackage{hyperref}                  % Para que funcionen los enlaces?
  \usepackage{enumitem}                  % Poner los numeros de los enumarate en negrita
\usepackage[normalem]{ulem}            % Para que funcione el tachado de texto en titulos con hyperref
\usepackage{soul}                      % Para el tachado de texto con \st{texto}
\usepackage{bm}                        % El comando \bm

% Enlaces, figuras y tablas subrayas y de color azul
\usepackage{xcolor}
\usepackage[normalem]{ulem} % para subrayado
\usepackage[colorlinks=true, linkcolor=blue, urlcolor=blue, citecolor=blue]{hyperref}
\renewcommand{\UrlFont}{\color{blue}\underline} % subraya links externos
\makeatletter % Subraya también referencias internas (figuras, secciones, etc.)
\AtBeginDocument{%
  \hypersetup{
    linkbordercolor = {0 0 1}, % borde azul si querés usar boxes
    linkcolor = blue,
  }%
  \let\oldhyper@link\hyper@link
  \def\hyper@link#1#2#3{%
  \oldhyper@link{#1}{#2}{\textcolor{blue}{\underline{#3}}}%
  }%
}
\makeatother

%sangria
\usepackage{etoolbox}
\makeatletter
\patchcmd{\@afterheading}{\@nobreaktrue}{}{}{}
\makeatother

% simbolos matematicos
\usepackage{amsmath, amssymb, amsthm, mathrsfs, amsfonts, amsbsy}

% Modifica el margen de la página. Mas granularmente seria [left=1in,right=1in,top=1in,bottom=1in] 
\usepackage[margin=2cm]{geometry}     

% encabezado
\usepackage{lastpage}
\usepackage{refcount} % Para obtener el valor de la pag sin q la pinte de azul
\usepackage{fancyhdr}
\pagestyle{fancy}
\fancyhf{}
\lhead{UNR - FCEyE}
% \rhead{\nouppercase{\leftmark}} % nombre de la seccion actual en minuscula, ya lo puse en chunks en el .qmd
\cfoot{Pág. \thepage\ - \getpagerefnumber{LastPage}}
\renewcommand{\footrulewidth}{0.08pt} %para que me haga la linea del footer
%\lfoot{}
%\rfoot{}

% para q enumere figuras segun seccion
%\counterwithin{figure}{section}


% colores portada
\usepackage{xcolor}
\usepackage{tikz}  %poner el coso de color
\definecolor{celestepastel}{RGB}{221, 245, 255}
\definecolor{rosabebe}{RGB}{255, 182, 193}
\definecolor{lila}{RGB}{208, 171, 198}


\author{}
\date{}
                          % Carga el partial con el titulo y sus configuraciones


\providecommand{\tightlist}{%
  \setlength{\itemsep}{0pt}\setlength{\parskip}{0pt}}\usepackage{longtable,booktabs,array}
\usepackage{calc} % for calculating minipage widths
% Correct order of tables after \paragraph or \subparagraph
\usepackage{etoolbox}
\makeatletter
\patchcmd\longtable{\par}{\if@noskipsec\mbox{}\fi\par}{}{}
\makeatother
% Allow footnotes in longtable head/foot
\IfFileExists{footnotehyper.sty}{\usepackage{footnotehyper}}{\usepackage{footnote}}
\makesavenoteenv{longtable}
\usepackage{graphicx}
\makeatletter
\newsavebox\pandoc@box
\newcommand*\pandocbounded[1]{% scales image to fit in text height/width
  \sbox\pandoc@box{#1}%
  \Gscale@div\@tempa{\textheight}{\dimexpr\ht\pandoc@box+\dp\pandoc@box\relax}%
  \Gscale@div\@tempb{\linewidth}{\wd\pandoc@box}%
  \ifdim\@tempb\p@<\@tempa\p@\let\@tempa\@tempb\fi% select the smaller of both
  \ifdim\@tempa\p@<\p@\scalebox{\@tempa}{\usebox\pandoc@box}%
  \else\usebox{\pandoc@box}%
  \fi%
}
% Set default figure placement to htbp
\def\fps@figure{htbp}
\makeatother

\makeatletter
\@ifpackageloaded{caption}{}{\usepackage{caption}}
\AtBeginDocument{%
\ifdefined\contentsname
  \renewcommand*\contentsname{Tabla de contenidos}
\else
  \newcommand\contentsname{Tabla de contenidos}
\fi
\ifdefined\listfigurename
  \renewcommand*\listfigurename{Listado de Figuras}
\else
  \newcommand\listfigurename{Listado de Figuras}
\fi
\ifdefined\listtablename
  \renewcommand*\listtablename{Listado de Tablas}
\else
  \newcommand\listtablename{Listado de Tablas}
\fi
\ifdefined\figurename
  \renewcommand*\figurename{Figura}
\else
  \newcommand\figurename{Figura}
\fi
\ifdefined\tablename
  \renewcommand*\tablename{Tabla}
\else
  \newcommand\tablename{Tabla}
\fi
}
\@ifpackageloaded{float}{}{\usepackage{float}}
\floatstyle{ruled}
\@ifundefined{c@chapter}{\newfloat{codelisting}{h}{lop}}{\newfloat{codelisting}{h}{lop}[chapter]}
\floatname{codelisting}{Listado}
\newcommand*\listoflistings{\listof{codelisting}{Listado de Listados}}
\makeatother
\makeatletter
\makeatother
\makeatletter
\@ifpackageloaded{caption}{}{\usepackage{caption}}
\@ifpackageloaded{subcaption}{}{\usepackage{subcaption}}
\makeatother
                         % Cargar el partial llamado 'pandoc' (por defecto de Quarto)

\floatplacement{figure}{H}             % Para que las figuras aparezcan en el lugar donde las ponemos 
% \setcounter{secnumdepth}{-\maxdimen}   % Para que las secciones no vayan numeradas. SECCIONES NUMERDAS porq comente la linea q desactiva la enumeracion
\setcounter{secnumdepth}{3}            % Para q enumere las sections(#) subsections(##) y subsubsections(###)
  \setlist[enumerate]{font=\bfseries}    % Poner los numeros de los enumarate en negrita
  \pdfstringdefDisableCommands{          % Para que funcione el tachado de texto en titulos con hyperref
    \renewcommand{\sout}{}
  }
  
  % graficos enumerados segun seccion en la q estoy
  \usepackage{chngcntr}
  \counterwithin{figure}{section}
  
  
\begin{document}

  \pagestyle{fancy}
  \rhead{}

  \tableofcontents

  \newpage
  \rhead{\nouppercase{\leftmark}}

  \section{Introducción}\label{introducciuxf3n}

  Este informe es un resumen de lo realizado en la Práctica
  Profesional\ldots{}

  El Centro de Agronegocios y Alimentos de la Sede Rosario de la
  Universidad Austral es una entidad especializada y de alto nivel
  académico cuya misión es impulsar, desde la región, el fortalecimiento
  de las cadenas de valor agroalimentarias e instituciones relacionadas,
  como un motor de desarrollo humano, económico y social de largo plazo
  para la Argentina, a través de la formación y de la investigación
  aplicada. La iniciativa promueve, desde un ámbito universitario, la
  capacidad de los sectores agroindustriales de crear y liderar las
  oportunidades y tendencias relacionadas con los agronegocios y
  alimentos, desde nuestro país, para la región y el mundo.

  En este centro se elabora, desde octubre de 2018, un índice llamado
  ``Ag Barometer''. El mismo proporciona bimensualmente medidas
  nacionales sobre las percepciones del sector agropecuario en la
  Argentina que muestran la evolución de su confianza a lo largo del
  tiempo. También se construye este índice en la Universidad de Purdue
  de Estados Unidos proporcionando un valor por mes desde 2015. Ambas
  universidades, con la misma metodología y de forma paralela, también
  construyen los índices ``Condiciones presentes'' y ``Expectativas
  futuras''.

  El Ag Barometer se construye a partir de una encuesta que se hace en
  cada edición a productores agropecuarios argentinos cuyo valor bruto
  de producción supera los USD 200.000 anuales. La muestra, de 406
  productores, se estratifica por escala productiva y garantiza
  representación mínima de las principales actividades del sector: soja,
  maíz, ganadería bovina, trigo, tambo y girasol.

  Se calcula a partir de cinco preguntas fijas, cada una de las cuales
  se puntúa según la diferencia entre respuestas favorables y
  desfavorables, escalada en un rango de 0 a 200, donde 100 representa
  el punto neutro. El valor final del índice corresponde al promedio
  simple de las cinco puntuaciones individuales, de modo que valores
  superiores a 100 reflejan percepciones positivas netas y valores
  inferiores a 100, percepciones negativas netas. Con la misma
  metodología y solo considerando algunas de las cinco preguntas, se
  calculan los índices Condiciones presentes y Expectativas futuras.

  El objetivo general de este trabajo es comparar con qué variables está
  relacionado el Ag Barometer de Argentina vs.~el de Estados Unidos.

  \newpage

  \section{Objetivos}\label{objetivos}

  \subsection{Objetivo general}\label{objetivo-general}

  El objetivo principal de esta Práctica Profesional es estudiar con qué
  variables está relacionado el Ag Barometer tanto en Argentina como en
  Estados Unidos y comparar qué ocurre en cada país.

  \subsection{Tareas específicas de la Práctica
  Profesional}\label{tareas-especuxedficas-de-la-pruxe1ctica-profesional}

  \begin{itemize}
  \tightlist
  \item
    Realizar un análisis descriptivo de los índices con el fin de
    conocer sus características.
  \item
    Imputación de valores donde es necesario.
  \item
    Buscar otras series que podrían estar correlacionadas con el Ag
    Barometer tanto en Argentina como en Estados Unidos.
  \item
  \end{itemize}

  \vdots

  \newpage

  \section{Materiales}\label{materiales}

  \subsection{Características del
  relevamiento}\label{caracteruxedsticas-del-relevamiento}

  La definición del universo a relevar tuvo como criterio la importancia
  de cada actividad en el producto bruto agropecuario. El mínimo de
  representación para cada actividad es el siguiente:

  \begin{itemize}
  \tightlist
  \item
    Soja 45 \%
  \item
    Maíz 21\%
  \item
    Ganadería de carne 16\%
  \item
    Trigo 7 \%
  \item
    Tambo 7\%
  \item
    Girasol 4 \%
  \end{itemize}

  Se realizó un muestreo estratificado de productores teniendo en cuenta
  datos del INDEC y \emph{Map of Ag}. Se detallan los estratos de
  productores en cuanto a su producción:

  \begin{itemize}
  \tightlist
  \item
    50\% de los encuestados con un valor bruto de producción entre USD
    200.000 y USD 399.000 al año
  \item
    30\% de los encuestados con valor bruto de producción entre USD
    414.000 y USD 999.999 al año
  \item
    20\% de los encuestados con un valor bruto de producción mayor o
    igual a USD 1.000.000 al año.
  \end{itemize}

  La encuesta a productores toma respuestas de 406 productores cuyo
  valor bruto de producción es igual o mayor a USD 200.000.

  \subsection{Conjunto de datos}\label{conjunto-de-datos}

  Los tres índices de Estados Unidos brindados por la Universidad de
  Purdue tienen periodicidad mensual y están completos. Comenzaron a
  construirlos en octubre de 2015 y la observación más reciente que se
  dispone es de marzo 2026.

  En cambio, los tres índices de Argentina provistos por el Centro de
  Agronegocios y Alimentos de la Sede Rosario de la Universidad Austral
  son medidos en los meses de enero, marzo, mayo, julio, septiembre y
  noviembre (impares); en 2020 en vez de haber observaciones para los
  meses de marzo y mayo, hay una para abril y se tienen datos desde
  octubre de 2018 hasta enero de 2026.

  A pesar de los meses en los que los índices locales son medidos, la
  primera observación de cada uno de ellos es en octubre de 2018 porque
  en este fecha se ha decidido comenzar el relevamiento y a partir del
  siguiente mes comenzaron a reportar los valores bimensualmente.

  Teniendo estas características en cuenta, los índices de Argentina
  tienen periodicidad irregular. Hay observaciones separadas por un mes
  (octubre de 2018 y noviembre de 2018), observaciones separadas por
  tres meses (enero 2020 y abril 2020, abril 2020 y julio 2020) y el
  resto de los datos están separados por dos meses.

  \newpage

  \section{Metodología}\label{metodologuxeda}

  Este trabajo tiene dos aspectos metodológicos.

  \subsection{Cálculo de los
  índices}\label{cuxe1lculo-de-los-uxedndices}

  \href{https://www.austral.edu.ar/rosario/agronegocios/investigacion-y-transferencia/barometer/}{Cálculo}

  \subsection{Definiciones básicas de una serie de
  tiempo}\label{definiciones-buxe1sicas-de-una-serie-de-tiempo}

  \vdots

  \newpage

  \section{Análisis descriptivo}\label{anuxe1lisis-descriptivo}

  En primer lugar, se realizó un análisis que permite explorar las
  características de los índices en los dos países.

  \subsection{Argentina}\label{argentina}

  Estas series tienen periodicidad irregular y valores faltantes. Voy a
  imputar porque\ldots{}

  Muestro características antes de imputar, dadas esas características y
  según las conclusiones de la tesina de Lucía Tamplín\ldots, imputo las
  series con interpolación lineal.

  Descriptivos de cada una de las 3 series imputadas

  CCFs

  \subsection{Estados Unidos}\label{estados-unidos}

  Descriptivo de cada uno de los 3 índices

  CCFs

  \newpage

  \section{Relación del Ag Barometer con otras
  variables}\label{relaciuxf3n-del-ag-barometer-con-otras-variables}

  \subsection{Variables numéricas}\label{variables-numuxe9ricas}

  Presentación de las variables (commodities, dow jones, fed, ICC\ldots)
  que cruzamos con el Ag Barometer

  Explicar por qué t vs.~t-1

  Muestro gráficos de correlaciones de Argentina y de EEUU

  \subsection{Variables categóricas}\label{variables-categuxf3ricas}

  Como no se observa relación entre el Ag Barometer y todas las
  variables propuestas, se sospecha que sí está relacionado con hitos
  políticos o bélicos.

  Para analizar esto se propone construir variables indicadoras que
  representan \ldots{}

  graficos

  \newpage

  \section{Comparación}\label{comparaciuxf3n}

  Comparar los países teniendo en cuenta las tres secciones anteriores.
  

\end{document}
